% ---------------------------------------------------------------------------
% Shared preamble for every edition.  Which edition is being built is decided
% by the entry point in tex/editions/, which defines exactly one of
% \SCAFFOLDResultsView, \SCAFFOLDRoutesView, \SCAFFOLDArchiveView before
% inputting tex/document.tex.
% ---------------------------------------------------------------------------
\documentclass[12pt]{article}

% Labels stay visible in the two working editions, because console and chat
% conversation about the project refers to LaTeX keys.  The results edition
% omits them for publication-style reading.
\ifdefined\SCAFFOLDShowKeys
  \usepackage[notcite,notref]{showkeys}
  % Long descriptive keys otherwise run past the physical page edge.
  \renewcommand*\showkeyslabelformat[1]{%
    \fbox{\normalfont\scriptsize\ttfamily#1}}
\fi

\usepackage{amsmath}
\usepackage{amsfonts}
\usepackage{amssymb}
\usepackage{amsthm}
\usepackage{array}
\usepackage{longtable}
\ifdefined\SCAFFOLDRoutesView
  \usepackage{xr-hyper}
\fi
% xcolor before hyperref: \termx uses the colour name `.', meaning the current
% colour, so a subsequent-occurrence link inherits the surrounding text.
\usepackage{xcolor}
\usepackage[colorlinks=true,linkcolor=blue,citecolor=blue,urlcolor=blue]{hyperref}

% --- project notation ------------------------------------------------------
% Keep macros here rather than in the modules, so that a change of notation is
% one edit and every edition agrees.
\newcommand{\N}{\mathbb{N}}
\newcommand{\R}{\mathbb{R}}
\newcommand{\Prset}{\mathbb{P}}          % the primes
\newcommand{\gc}{g}                      % the Goldbach count
% ---------------------------------------------------------------------------

\newtheorem{theorem}{Theorem}[section]
\newtheorem{lemma}[theorem]{Lemma}
\newtheorem{cor}[theorem]{Corollary}
\newtheorem{prop}[theorem]{Proposition}
\newtheorem{conjecture}[theorem]{Conjecture}
\newtheorem{prob}[theorem]{Problem}
\theoremstyle{definition}
\newtheorem{definition}[theorem]{Definition}
\newtheorem{example}[theorem]{Example}
\theoremstyle{plain}

\renewcommand{\baselinestretch}{1.15}
\emergencystretch=1em
\setcounter{tocdepth}{1}

% The routes companion refers to labels owned by the results edition, so the
% Makefile builds the results edition first.  Which results auxiliary to read is
% a parameter, because the plain and linked variants are built into different
% directories and each must import its own: a link in the linked routes PDF has
% to land in the linked results PDF.
\ifdefined\SCAFFOLDRoutesView
  \ifdefined\SCAFFOLDResultsAux\else
    \def\SCAFFOLDResultsAux{build/plain/results/scaffold-results}
  \fi
  % Labels are read from the build tree, but the link must point at the released
  % sibling: output/pdf/{plain,linked}/ holds scaffold-results.pdf next to
  % scaffold-routes.pdf.  Without the explicit URL, xr-hyper defaults the target
  % to the build path, which does not exist relative to the released PDF, and
  % every cross-document link silently goes nowhere.
  \def\SCAFFOLDimportresults#1{\externaldocument{#1}[scaffold-results.pdf]}
  \expandafter\SCAFFOLDimportresults\expandafter{\SCAFFOLDResultsAux}
\fi

% ---------------------------------------------------------------------------
% Terminology apparatus.  A term is marked \dfn{goldbach pair} at the single
% place where it is defined, and \term{goldbach pair} at the first use in every
% other module; the latter links to the term's row in tex/glossary.tex, which
% points on to the defining section or result.  Because the glossary is included
% by all three editions, \term always resolves, even when the definition is
% owned by another edition and reached through xr-hyper.  Use \dfnas/\termas
% when the printed form differs from the key (plurals, symbols, inflections).
% ---------------------------------------------------------------------------
\newcommand{\dfnas}[2]{\hypertarget{term:#1}{}\emph{#2}}
\newcommand{\dfn}[1]{\dfnas{#1}{#1}}
\newcommand{\termas}[2]{\hyperlink{gloss:#1}{#2}}
\newcommand{\term}[1]{\termas{#1}{#1}}

% \termx marks a *subsequent* occurrence of a term.  It is a no-op unless the
% edition defines \SCAFFOLDLinkAll, and even then it renders the word in the
% running text colour rather than in blue: the link is there for a reader who
% wants it, without turning one word in twenty into a coloured link.
% Occurrences are written by tools/python/link_all.py at build time, never by
% hand.
\ifdefined\SCAFFOLDLinkAll
  \newcommand{\termx}[2]{\begingroup\hypersetup{linkcolor=.}%
    \hyperlink{gloss:#1}{#2}\endgroup}
\else
  \newcommand{\termx}[2]{#2}
\fi

% First cell of a glossary row: the anchor that \term jumps to.  \gkey marks a
% term that the body defines, and whose key must be matched by exactly one
% \dfn; \gkeyx marks one the glossary itself defines, for the synonyms, method
% names and imported notions the text uses without ever stopping to introduce
% them.
\newcommand{\gkey}[2]{\hspace{0pt}\textbf{#2}\hypertarget{gloss:#1}{}}
\newcommand{\gkeyx}[2]{\gkey{#1}{#2}}
